\section{USAMO 1990}
        \begin{enumerate}
        	\item (\textit{USAMO 1990}) A certain state issues license plates consisting of six digits (from 0 through 9). The state requires that any two plates differ in at least two places. (Thus the plates $\boxed{027592}$ and $\boxed{020592}$ cannot both be used.) Determine, with proof, the maximum number of distinct license plates that the state can use.


 



\textbf{Solution 1}

Consider license plates of $n$ digits, for some fixed $n$, issued with the same criteria.

 We first note that by the pigeonhole principle, we may have at most $10^{n-1}$ distinct plates.  Indeed, if we have more, then there must be two plates which agree on the first $n-1$ digits; these plates thus differ only on one digit, the last one.

 We now show that it is possible to issue $10^{n-1}$ distinct license plates which satisfy the problem's criteria.  Indeed, we issue plates with all $10^{n-1}$ possible combinations for the first $n-1$ digit, and for each plate, we let the last digit be the sum of the preceding digits taken mod 10.  This way, if two plates agree on the first $n-1$ digits, they agree on the last digit and are thus the same plate, and if two plates differ in only one of the first $n-1$ digits, they must differ as well in the last digit.

 It then follows that $10^{n-1}$ is the greatest number of license plates the state can issue.  For $n=6$, as in the problem, this number is $10^5$.  $\blacksquare$

 



\textbf{Solution 2}

$\text{Constructive Induction for } n=6$

 \begin{verbatim}

 \[n=1 \quad \text{1 element} \quad n=2 \quad \text{10 elements}\]  Suppose for $n$, we have a set $\left\{ \left( a_1, \ldots, a_n \right) \right\}$ that satisfies the condition. We need to construct $S_{n+1}$ such that it does not contain $\left( a_1, \ldots, a_n \right)$, denoted as $\left( b_1, \ldots, b_n \right)$.  Construct:  
 \[S_{n+1} = \left\{ \left( a_1, a_2, \ldots, a_{n-k}, k \right) \mid k = 1, 2, \ldots, 10, \left( a_1, \ldots, a_n \right) \in S_n \right\}\]  If $\left( a_1 + k, \ldots, a_{n+k}, k \right)$ and $\left( b_1, \ldots, b_{n+\ell}, \ell \right)$ have exactly one difference:  1. Case 1: $k \neq \ell$     
 \[a_1 = b_1, \ldots, a_{n-1} = b_{n-1}, a_n = b_n \Rightarrow a_{n+k} \neq b_{n+\ell}, \text{contradiction}\]  2. Case 2: $k = \ell$     If $a_i = b_i$ for $i \neq n$,     
 \[a_1 = b_1, \ldots, a_n = b_n \quad (\text{removing } i)\]     Contradiction! If $a_{n+k} \neq b_{n+\ell} \Rightarrow a_n \neq b_n$, contradiction with the same $k$.  At least two positions are different, i.e., there cannot be exactly one difference. Therefore, in the set $\{0, 1, \ldots, 9\}^6$ (hypercube), there are no collinear points (parallel to the coordinate axes).  
 \[\text{1 dimension: 1 element}\]  
 \[\text{2 dimensions: 10 elements}\]  
 \[\text{3 dimensions: 100 elements}\]  
 \[10^6 \text{ points can be covered by } 10^5 \text{ lines} \Rightarrow \text{At most } 10^5 \text{ elements.}\]
\end{verbatim}




	\item (\textit{USAMO 1990}) A sequence of functions $\, \{f_n(x) \} \,$ is defined recursively as follows:

 \begin{align*} f_1(x) &= \sqrt {x^2 + 48}, \quad \text{and} \\ f_{n + 1}(x) &= \sqrt {x^2 + 6f_n(x)} \quad \text{for } n \geq 1. \end{align*}

(Recall that $\sqrt {\makebox[5mm]{}}$ is understood to represent the positive square root.) For each positive integer $n$, find all real solutions of the equation $\, f_n(x) = 2x \,$.


 



\textbf{Solution 1}

We define $f_0(x) = 8$.  Then the recursive relation holds for $n=0$, as well.

 Since $f_n (x) \ge 0$ for all nonnegative integers $n$, it suffices to consider nonnegative values of $x$.

 We claim that the following set of relations hold true for all natural numbers $n$ and nonnegative reals $x$:
 \begin{align*} f_n(x) &< 2x \text{ if }x>4 ; \\ f_n(x) &= 2x \text{ if }x=4 ; \\ f_n(x) &> 2x \text{ if }x<4 . \end{align*}
To prove this claim, we induct on $n$.  The statement evidently holds for our base case, $n=0$.

 Now, suppose the claim holds for $n$.  Then
 \begin{align*} f_{n+1}(x) &= \sqrt{x^2 + 6f_n(x)} < \sqrt{x^2+12x} < \sqrt{4x^2} = 2x, \text{ if } x>4 ; \\ f_{n+1}(x) &= \sqrt{x^2 + 6f_n(x)} = \sqrt{x^2 + 12x} = \sqrt{4x^2} = 2x, \text{ if } x=4 ; \\ f_{n+1}(x) &= \sqrt{x^2 + 6f_n(x)} > \sqrt{x^2+12x} > \sqrt{4x^2} = 2x, \text{ if } x<4 . \end{align*}
The claim therefore holds by induction.  It then follows that for all nonnegative integers $n$, $x=4$ is the unique solution to the equation $f_n(x) = 2x$.  $\blacksquare$

 



\textbf{Solution 2}

We claim that the only solution is $\boxed{x=4}$ for all such $n$. To show this, we consider all $x$ and find solutions $n$.

 Before we consider solutions, we show that for all $x$, $f_n(x)$ is positive for all positive integers $n$ by induction. For our base case:
 \[f_1(x)=\sqrt{x^2+48}\ge\sqrt{48}>0\]so it is positive. Next, for our inductive step, assume for some $n=k-1$ that $f_{k-1}(x)$ is positive; thus $f_{k-1}(x)>0$, and we show that $f_k(x)$ is positive:
 \[f_k(x)=\sqrt{x^2+6f_{k-1}(x)}\ge\sqrt{0+6f_{k-1}(x)}>\sqrt{0+0}=0\]thus $f_k(x)>0$, so we have proven the claim by induction.

 First, we consider negative $x$. We know that for negative $x$, if the equation were to have a solution, we would have $f_n(x)=2x<0$ for some $n$. However, $f_n(x)$ is always nonnegative since $f_n(x)$ is always the square root of some number, which can never be negative; thus there are no solutions in this case.

 Next, we consider $x=0$. Solutions would have to satisfy $f_n(0)=0$; we have previously shown that all $f_n(x)$ are positive, so there are no such $n$.

 Finally, we consider positive $x$. We divide this set into three groups: $x<4$, $x=4$, and $x>4$. We claim that for increasing $n$, $f_n(x)$ is decreasing, constant, and increasing, respectively, and we prove this using induction. First, we analyze $x<4$. Then
 \[f_1(x)<\sqrt{4^2+48}=\sqrt{64}=8\]For our base case:
 \[f_2(x)=\sqrt{x^2+6f_1(x)}<\sqrt{x^2+6\cdot8}=f_1(x)\]Thus $f_2(x)<f_1(x)$. For our inductive step, if for some $n=k-1$ we have $f_k(x)<f_{k-1}(x)$, we show that $f_{k+1}(x)<f_k(x)$:
 \[f_{k+1}(x)=\sqrt{x^2+6f_k(x)}<\sqrt{x^2+6f_{k-1}(x)}=f_k(x)\]Thus $f_1(x)>f_2(x)>f_3(x)>\ldots$ and the conclusion follows.

 Next, for $x=4$, $f_1(4)=8$ by substituting. Then, if $f_{n-1}(4)=8$, we show that $f_n(4)=8$ as well:$f_n(4)=\sqrt{x^2+6f_{n-1}(x)}=\sqrt{16+6\cdot8}=\sqrt{64}=8$so the sequence is constant. Notice that in this case, the equation we must solve becomes $f_n(4)=8$, which is true for all positive integers $n$.

 Finally, we analyze $x>4$. Then
 \[f_1(x)>\sqrt{4^2+48}=\sqrt{64}=8\]For our base case:
 \[f_2(x)=\sqrt{x^2+6f_1(x)}>\sqrt{x^2+6\cdot8}=f_1(x)\]Thus $f_2(x)>f_1(x)$. For our inductive step, if for some $n=k-1$ we have $f_k(x)>f_{k-1}(x)$, we show that $f_{k+1}(x)>f_k(x)$:
 \[f_{k+1}(x)=\sqrt{x^2+6f_k(x)}>\sqrt{x^2+6f_{k-1}(x)}=f_k(x)\]Thus $f_1(x)<f_2(x)<f_3(x)<\ldots$ and the conclusion follows.

 If for some positive integer $n$ we have $f_n(x)=2x$, then:

 \begin{align*}f_n(x)&=\sqrt{x^2+6f_{n-1}(x)}\\2x&=\sqrt{x^2+6f_{n-1}(x)}\\4x^2&=x^2+6f_{n-1}(x)\\3x^2&=6f_{n-1}(x)\\\frac{1}{2}x^2&=f_{n-1}(x)\end{align*}
However, if $0<x<4$ is a solution, then we must have $f_{n-1}(x)>f_n(x)$; substituting values yields $\frac{1}{2}x^2>2x$, which implies $x>4$; this is a contradiction. Similarly, if $x>4$ is a solution, then we must have $f_{n-1}(x)<f_n(x)$; substituting values yields $\frac{1}{2}x^2<2x$, which implies $x<4$; this is also a contradiction.

 As a result, since we have already established that $x=4$ is a solution for all $n$ and no other $x$ can be solutions, we conclude that for each positive integer $n$, the only real solution of the equation is $x=4$.

 


	\item (\textit{USAMO 1990}) Suppose that necklace $\, A \,$ has 14 beads and necklace $\, B \,$ has 19. Prove that for any odd integer $n \geq 1$, there is a way to number each of the 33 beads with an integer from the sequence

 \[\{ n, n + 1, n + 2, \dots, n + 32 \}\]
so that each integer is used once, and adjacent beads correspond to relatively prime integers. (Here a "necklace" is viewed as a circle in which each bead is adjacent to two other beads.)


 



\textbf{Solution}

\textbf{Lemma.}  For every positive odd integer $n$, there exists a nonnegative integer $k \le 17$ such that $n+k$ is relatively prime to $n+k+15$, $n+k+1$ is relatively prime to $n+k+14$.

 \textit{Proof.}  Consider the positive integers $n, n+1, n+2$.  Note that at most one of these is divisible by 3, and at most one is divisible by 5.  Therefore one of these, say $n+a$, is divisible by neither, and is therefore relatively prime to 15.  Furthermore, $n+a+1$ and $n+(a+15)+1$ have different residues mod 13, so one of them, say $n+b+1$, is relatively prime to 13.  Since $n+b \equiv n+a \pmod{15}$, $n+b$ is relatively prime to 15.  But
 \[\gcd(n+b,n+b+15) = \gcd\bigl[ n+b,(n+b+15) - (n+b) \bigr] = \gcd(n+b,15) = 1,\]so $n+b$ is relatively prime to $n+b+15$.  Also,
 \[\gcd(n+b+1, n+b+14) = \gcd \bigl[ n+b+1, (n+b+14)-(n+b+1) \bigr] = \gcd( n+b+1, 13) ,\]so $n+b+1$ is relatively prime to $n+b+14$.  Finally, $b \le a+15 \le 17$.  It follows that setting $k=b$ satisfies the lemma.  $\blacksquare$

 Let $k$ be an integer as described in the lemma.  We place the integers $n, \dotsc, n+k, n+k+15, \dotsc, n+32$ on the necklace with 19 beads, and the integers $n+k+1, \dotsc, n+k+14$ on the necklace with 14 beads, in those orders.  Since $n$ is odd, $n$ is relatively prime to $n+32 = n+2^5$.  By definition, $n+k$ and $n+k+15$ are relatively prime, as are $n+k+1$ and $n+k+14$; finally, $a$ and $a+1$ are relatively prime for all integers $a$.  It follows that each bead in this arrangement is relatively prime to its neighbors.  $\blacksquare$

 
<i>Alternate solutions are always welcome. If you have a different, elegant solution to this problem, please \href{https://artofproblemsolving.com/wiki/index.php?title=1990\_USAMO\_Problems/Problem\_3&action=edit}{add it} to this page.</i>

 


	\item (\textit{USAMO 1990}) Find, with proof, the number of positive integers whose base-$n$ representation consists of distinct digits with the property that, except for the leftmost digit, every digit differs by $\pm 1$ from some digit further to the left. (Your answer should be an explicit function of $n$ in simplest form.)


 



\textbf{Solution}

Let a $k$-good sequence be a sequence of distinct integers $\{ a_i \}_{i=1}^k$ such that for all integers $2\le i \le k$, $a_i$ differs from some preceding term by $\pm 1$.

 \textbf{Lemma.}  Let $a$ be an integer.  Then there are $2^{k-1}$ $k$-good sequences starting on $a$, and furthermore, the terms of each of these sequences constitute a permutation of $k$ consecutive integers.

 \textit{Proof.} We induct on $k$.  For $k=1$, the lemma is trivially true.  Now, suppose the lemma holds for $k$.  If $\{ a_i \}_{i=1}^{k+1}$ is a $(k+1)$-good sequence, then $\{ a_i \}_{i=1}^k$ is a $k$-good sequence which starts on $a$, so it is a permutation of $k$ consecutive integers, say $m, \dotsc, M$.  Then the only possibilities for $a_{k+1}$ are $m-1$ and $M+1$; either way, $\{ a_i \}_{i=1}^{k+1}$ constitutes a permutation of $k+1$ consecutive integers.  Since there are $2^k$ possible sequences $\{a_i\}_{i=1}^k$, and 2 choices of $a_{k+1}$ for each of these sequences, it also follows that there are $2^k \cdot 2 = 2^{k+1}$ $(k+1)$-good sequences which start on $a$.  Thus the lemma holds by induction.  $\blacksquare$

 We now consider the number of desired positive integers with $k$ digits.  Evidently, $k$ must be less than or equal to $n$.  We also note that the digits of such an integer must constitute a $k$-good sequence.  Since the minimum of this sequence can be any of the digits $0, \dotsc, n-k$, unless the minimum is 0 and is the first digit (in which case the only possible sequence is an increasing arithmetic sequence), and there are $2^{k-1}$ $k$-good sequences up to translation, it follows that there are $(n-k+1) 2^{k-1}-1$ desired positive integers with $k$ digits.  Thus the total number of desired positive integers is
 \begin{align*} \sum_{k=1}^n \bigl[ (n-k+1) 2^{k-1}-1 \bigr] &= -n + \sum_{k=1}^n \sum_{j=k}^n 2^{k-1} = -n + \sum_{j=1}^n \sum_{k=1}^j 2^{k-1} \\ &= -n + \sum_{j=1}^n (2^k-1) = - 2n -1 + \sum_{j=0}^n 2^k, \end{align*}
which is equal to $2^{n+1} - 2(n+1)$, our answer.  $\blacksquare$

 


 The lemma in the above solution is wrong. A 3-term sequence starting on 1 will not comply, for instance.My solution is shown as follows.We call all qualified numbers ‘good numbers’, with k digits, 1<=k<=n. See any good number as a k-term sequence, m being the minimal term, a being the first term. Note that 0<=m<=n-k.Lemma. It is trivially true that any good number consists of consecutive digits. It can be proven with induction easily.Now pick any m, let f(i) be the number of sequences when a=i, we prove that f(i)=C(k-1,i-m), m<=i<=m+k-1.Per Lemma, a sequence consists of consecutive digits: m,m+1,\ldots,i,\ldots,m+k-1. The first term a=i determined, there are k-1 remaining spaces to be filled. The digits smaller than i are m,m+1,\ldots,i-1, their only link to i is i-1, then i-2, so on so forth, eventually we can know they can only enter the sequence in one order which is descending: i-1,i-2,\ldots,m. Likewise, the digits bigger than i can only enter the sequence in one order which is ascending: i+1,i+2,\ldots,m+k-1. Therefore, the number of sequences is equivalent to the number of ways in which the i-m digits smaller than i occupy the k-1 remaining spaces, i.e., C(k-1,i-m).So, for any given m, the total number of sequences is $\sum$(f(i),i=m~m+k-1)=$\sum$(C(k-1,i),i=0~k-1)=2\textasciicircum{}(k-1). Considering that there are n-k+1 possibilities for m, and deducting the one special case with m=0 and a=0 (because the first digit cannot be 0), we now can have the total number of sequences for any given k: (n-k+1)2\textasciicircum{}(k-1)-1. Now we can arrive at the grand total: $\sum$((n-k+1)2\textasciicircum{}(k-1)-1,k=1~n). It is trivial that the final sum is 2\textasciicircum{}(n+1)-2(n+1). Done.Tmath2024 2026.5.18

 
<i>Alternate solutions are always welcome. If you have a different, elegant solution to this problem, please \href{https://artofproblemsolving.com/wiki/index.php?title=1990\_USAMO\_Problems/Problem\_4&action=edit}{add it} to this page.</i>

 


	\item (\textit{USAMO 1990}) An acute-angled triangle $ABC$ is given in the plane. The circle with diameter $\, AB \,$ intersects altitude $\, CC' \,$ and its extension at points $\, M \,$ and $\, N \,$, and the circle with diameter $\, AC \,$ intersects altitude $\, BB' \,$ and its extensions at $\, P \,$ and $\, Q \,$. Prove that the points $\, M, N, P, Q \,$ lie on a common circle.


 



\textbf{Solution 1}

Let $A'$ be the intersection of the two circles (other than $A$). $AA'$ is perpendicular to both $BA'$, $CA'$ implying $B$, $C$, $A'$ are collinear. Since $A'$ is the foot of the altitude from $A$: $A$, $H$, $A'$ are concurrent, where $H$ is the orthocentre.

 Now, $H$ is also the intersection of $BB'$, $CC'$ which means that $AA'$, $MN$, $PQ$ are concurrent. Since $A$, $M$, $N$, $A'$ and $A$, $P$, $Q$, $A'$ are cyclic, $M$, $N$, $P$, $Q$ are cyclic by the radical axis theorem.

 



\textbf{Solution 2}

Define $A'$ as the foot of the altitude from $A$ to $BC$. Then, $AA' \cap BB' \cap CC'$ is the orthocenter. We will denote this point as $H$.Since $\angle AA'C$ and $\angle AA'B$ are both $90^{\circ}$, $A'$ lies on the circles with diameters $AC$ and $AB$.

 Now we use the Power of a Point theorem with respect to point $H$. From the circle with diameter $AB$ we get $AH \cdot A'H = MH \cdot NH$. From the circle with diameter $AC$ we get $AH \cdot A'H = PH \cdot QH$. Thus, we conclude that $PH \cdot QH = MH \cdot NH$, which implies that $P$, $Q$, $M$, and $N$ all lie on a circle.

 



\textbf{Solution 3 (Radical Lemma)}

Let $\omega_1$ be the circumcircle with diameter $AB$ and $\omega_2$ be the circumcircle with diameter $AC$.  We claim that the second intersection of $\omega_1$ and $\omega_2$ other than $A$ is $A'$, where $A'$ is the feet of the perpendicular from $A$ to segment $BC$.  Note that 
 \[\angle AA'B=90^{\circ}=\angle AB'B\] so $A'$ lies on $\omega_1.$  Similarly, $A'$ lies on $\omega_2$.  Hence, $AA'$ is the radical axis of $\omega_1$ and $\omega_2$.  By the Radical Lemma, it suffices to prove that the intersection of lines $MN$ and $PQ$ lie on $AA'$.  But, $MN$ is the same line as $CC'$ and $PQ$ is the same line as $BB'$.  Since $AA', BB'$, and $CC'$ intersect at the orthocenter $H$, $H$ lies on the radical axis $AA'$ and we are done. $\blacksquare$

 



\textbf{Solution 4}

We know that $BCB'C'$ is a cyclic quadrilateral. Hence, 

 $HB \cdot HB' = HC \cdot HC'$

 $\implies Pow_{\omega_{1}} = Pow_{\omega_{2}}$

 $\implies HM \cdot HN = HP \cdot HQ$

 $\implies MPNQ$ is cyclic $\raggedright\blacksquare$.

 


\end{enumerate}